\documentclass[10pt,a4paper]{amsart}
\usepackage{amsmath,amssymb,amsthm,graphicx}
% For dealing with references we use the comment environment
\usepackage{verbatim}
\newenvironment{reference}{\comment}{\endcomment}
%\newenvironment{reference}{}{}
\newenvironment{slogan}{\comment}{\endcomment}
\newenvironment{history}{\comment}{\endcomment}
% For commutative diagrams we use Xy-pic
\usepackage[all]{xy}

% We use 2cell for 2-commutative diagrams.
\xyoption{2cell}
\UseAllTwocells
\theoremstyle{plain}
\newtheorem{theorem}[subsection]{Theorem}
\newtheorem{proposition}[subsection]{Proposition}
\newtheorem{lemma}[subsection]{Lemma}

\theoremstyle{definition}
\newtheorem{definition}[subsection]{Definition}
\newtheorem{example}[subsection]{Example}
\newtheorem{exercise}[subsection]{Exercise}
\newtheorem{situation}[subsection]{Situation}

\theoremstyle{remark}
\newtheorem{remark}[subsection]{Remark}
\newtheorem{remarks}[subsection]{Remarks}

\numberwithin{equation}{subsection}
% Macros
%
\def\lim{\mathop{\mathrm{lim}}\nolimits}
\def\colim{\mathop{\mathrm{colim}}\nolimits}
\def\Spec{\mathop{\mathrm{Spec}}}
\def\Hom{\mathop{\mathrm{Hom}}\nolimits}
\def\Ext{\mathop{\mathrm{Ext}}\nolimits}
\def\SheafHom{\mathop{\mathcal{H}\!\mathit{om}}\nolimits}
\def\SheafExt{\mathop{\mathcal{E}\!\mathit{xt}}\nolimits}
\def\Sch{\mathit{Sch}}
\def\Mor{\mathop{\mathrm{Mor}}\nolimits}
\def\Ob{\mathop{\mathrm{Ob}}\nolimits}
\def\Sh{\mathop{\mathit{Sh}}\nolimits}
\def\NL{\mathop{N\!L}\nolimits}
\def\CH{\mathop{\mathrm{CH}}\nolimits}
\def\proetale{{pro\text{-}\acute{e}tale}}
\def\etale{{\acute{e}tale}}
\def\QCoh{\mathit{QCoh}}
\def\Ker{\mathop{\mathrm{Ker}}}
\def\Im{\mathop{\mathrm{Im}}}
\def\Coker{\mathop{\mathrm{Coker}}}
\def\Coim{\mathop{\mathrm{Coim}}}

% Boxtimes
%
\DeclareMathSymbol{\boxtimes}{\mathbin}{AMSa}{"02}

%
% Macros for moduli stacks/spaces
%
\def\QCohstack{\mathcal{QC}\!\mathit{oh}}
\def\Cohstack{\mathcal{C}\!\mathit{oh}}
\def\Spacesstack{\mathcal{S}\!\mathit{paces}}
\def\Quotfunctor{\mathrm{Quot}}
\def\Hilbfunctor{\mathrm{Hilb}}
\def\Curvesstack{\mathcal{C}\!\mathit{urves}}
\def\Polarizedstack{\mathcal{P}\!\mathit{olarized}}
\def\Complexesstack{\mathcal{C}\!\mathit{omplexes}}
% \Pic is the operator that assigns to X its picard group, usage \Pic(X)
% \Picardstack_{X/B} denotes the Picard stack of X over B
% \Picardfunctor_{X/B} denotes the Picard functor of X over B
\def\Pic{\mathop{\mathrm{Pic}}\nolimits}
\def\Picardstack{\mathcal{P}\!\mathit{ic}}
\def\Picardfunctor{\mathrm{Pic}}
\def\Deformationcategory{\mathcal{D}\!\mathit{ef}}
\usepackage{cleveref}
\crefname{originallemma}{original lemma}{original lemmas}
\Crefname{originallemma}{Original lemma}{Original lemmas}
\makeatletter
\newlabel{coherent-lemma-existence-easy}{{0889}{source}{}{sourcepointer-1}{}}
\newlabel{coherent-lemma-existence-easy@cref}{{[originallemma][0][]0889}{[1][0][]source}{}{}{}}
\makeatother
\crefname{originallemma}{original lemma}{original lemmas}
\Crefname{originallemma}{Original lemma}{Original lemmas}
\makeatletter
\newlabel{lemma-improvement-formal-coherent-module-better}{{0EJF}{source}{}{sourcepointer-2}{}}
\newlabel{lemma-improvement-formal-coherent-module-better@cref}{{[originallemma][0][]0EJF}{[1][0][]source}{}{}{}}
\makeatother
\crefname{originallemma}{original lemma}{original lemmas}
\Crefname{originallemma}{Original lemma}{Original lemmas}
\makeatletter
\newlabel{more-morphisms-lemma-apply-proj-spec}{{0EKI}{source}{}{sourcepointer-3}{}}
\newlabel{more-morphisms-lemma-apply-proj-spec@cref}{{[originallemma][0][]0EKI}{[1][0][]source}{}{}{}}
\makeatother
\crefname{originalsection}{original section}{original sections}
\Crefname{originalsection}{Original section}{Original sections}
\makeatletter
\newlabel{more-morphisms-section-proj-spec}{{0EKF}{source}{}{sourcepointer-4}{}}
\newlabel{more-morphisms-section-proj-spec@cref}{{[originalsection][0][]0EKF}{[1][0][]source}{}{}{}}
\makeatother
\crefname{originallemma}{original lemma}{original lemmas}
\Crefname{originallemma}{Original lemma}{Original lemmas}
\makeatletter
\newlabel{coherent-lemma-coherent-proper-ample}{{0B5T}{source}{}{sourcepointer-5}{}}
\newlabel{coherent-lemma-coherent-proper-ample@cref}{{[originallemma][0][]0B5T}{[1][0][]source}{}{}{}}
\makeatother
\crefname{originallemma}{original lemma}{original lemmas}
\Crefname{originallemma}{Original lemma}{Original lemmas}
\makeatletter
\newlabel{coherent-lemma-inverse-systems-pullback}{{0887}{source}{}{sourcepointer-6}{}}
\newlabel{coherent-lemma-inverse-systems-pullback@cref}{{[originallemma][0][]0887}{[1][0][]source}{}{}{}}
\makeatother
\crefname{originallemma}{original lemma}{original lemmas}
\Crefname{originallemma}{Original lemma}{Original lemmas}
\makeatletter
\newlabel{lemma-discussion}{{0EIY}{source}{}{sourcepointer-7}{}}
\newlabel{lemma-discussion@cref}{{[originallemma][0][]0EIY}{[1][0][]source}{}{}{}}
\makeatother
\crefname{originallemma}{original lemma}{original lemmas}
\Crefname{originallemma}{Original lemma}{Original lemmas}
\makeatletter
\newlabel{lemma-cd-1-canonical}{{0EJI}{source}{}{sourcepointer-8}{}}
\newlabel{lemma-cd-1-canonical@cref}{{[originallemma][0][]0EJI}{[1][0][]source}{}{}{}}
\makeatother
\crefname{originalequation}{original equation}{original equations}
\Crefname{originalequation}{Original equation}{Original equations}
\makeatletter
\newlabel{more-morphisms-equation-cohomology-torsor}{{0EKH}{source}{}{sourcepointer-9}{}}
\newlabel{more-morphisms-equation-cohomology-torsor@cref}{{[originalequation][0][]0EKH}{[1][0][]source}{}{}{}}
\makeatother
\crefname{originallemma}{original lemma}{original lemmas}
\Crefname{originallemma}{Original lemma}{Original lemmas}
\makeatletter
\newlabel{properties-lemma-quasi-coherent-quasi-affine}{{0EHM}{source}{}{sourcepointer-10}{}}
\newlabel{properties-lemma-quasi-coherent-quasi-affine@cref}{{[originallemma][0][]0EHM}{[1][0][]source}{}{}{}}
\makeatother
\crefname{originallemma}{original lemma}{original lemmas}
\Crefname{originallemma}{Original lemma}{Original lemmas}
\makeatletter
\newlabel{lemma-Gm-equivariant-extend-canonically}{{0EL4}{source}{}{sourcepointer-11}{}}
\newlabel{lemma-Gm-equivariant-extend-canonically@cref}{{[originallemma][0][]0EL4}{[1][0][]source}{}{}{}}
\makeatother
\crefname{originalproposition}{original proposition}{original propositions}
\Crefname{originalproposition}{Original proposition}{Original propositions}
\makeatletter
\newlabel{coherent-proposition-coherent-modules-on-proj-general}{{0BXF}{source}{}{sourcepointer-12}{}}
\newlabel{coherent-proposition-coherent-modules-on-proj-general@cref}{{[originalproposition][0][]0BXF}{[1][0][]source}{}{}{}}
\makeatother
\crefname{originallemma}{original lemma}{original lemmas}
\Crefname{originallemma}{Original lemma}{Original lemmas}
\makeatletter
\newlabel{lemma-recognize-formal-coherent-modules}{{0EIQ}{source}{}{sourcepointer-13}{}}
\newlabel{lemma-recognize-formal-coherent-modules@cref}{{[originallemma][0][]0EIQ}{[1][0][]source}{}{}{}}
\makeatother
\title{Proposition lefschetz existence in algebraization.tex}
\author{The Stacks Project Authors}
\date{stacks-project:a04446e57ec1fbc252a871afcec7752fb2807b14; GNU Free Documentation License 1.2 (as published in the Stacks Project COPYING)}
\makeatletter
\newlabel{divisors-definition-zero-scheme-s}{{02OQ}{source}{}{context-source-1}{}}
\makeatother
\makeatletter
\newlabel{coherent-section-coherent-formal}{{0EHN}{source}{}{context-source-2}{}}
\makeatother
\begin{document}
\maketitle
\noindent\textit{Editorial note.} All mathematical statement and proof text is reproduced verbatim from the author's chapter algebraization.tex of the Stacks Project, fixed at commit a04446e57ec1fbc252a871afcec7752fb2807b14. The pinned preamble is supplied as its own file; the wrapper class, the xr-hyper wiring and the \textbackslash{}externaldocument list are omitted because the assembler supplies the class.\par\medskip
\noindent\textit{Editorial note.} Numbering inside this standalone excerpt is editorial, not the full Stacks book print numbering.\par\medskip
\noindent\textit{Editorial note.} The candidate is from the extractor's stacks-candidates line, not from the earlier judged initial\_corpus line; the H2 judgement and its reason below were made in this module by reading the author's own statement and complete proof in the fixed checkout. Original candidate key: stacks:algebraization.tex:proposition-lefschetz-existence.\par\medskip
\noindent\textit{Editorial note.} Reference labels in the proof resolve to navigation pointers to the original author source; these pointers do not reproduce proofs of the prerequisites.\par\medskip
\noindent\textit{Allowed standard prerequisites (editorial).}\begin{itemize}
\item Prerequisites named by the author inside the counted span. The counted span is reproduced verbatim, including the author's own citations. Its prerequisites are the earlier results of the same fixed Stacks Project version that the author cites there and in the chapter's lead-in.
\end{itemize}\medskip
\ifdefined\hypertarget\hypertarget{sourcepointer-1}{}\fi
\noindent\textit{Source reference.} Original lemma 0889: sources/coherent.tex, source lines 6487--6487. Navigation pointer used in the counted span; the span contains the pointer only and the cited result is not reproduced inside it. Registered as a non-load-bearing navigation stub in the established form for Stacks material. No bibliography citation is involved.\par\medskip
\ifdefined\hypertarget\hypertarget{sourcepointer-2}{}\fi
\noindent\textit{Source reference.} Original lemma 0EJF: sources/algebraization.tex, source lines 6952--6952. Navigation pointer used in the counted span; the span contains the pointer only and the cited result is not reproduced inside it. Registered as a non-load-bearing navigation stub in the established form for Stacks material. No bibliography citation is involved.\par\medskip
\ifdefined\hypertarget\hypertarget{sourcepointer-3}{}\fi
\noindent\textit{Source reference.} Original lemma 0EKI: sources/more-morphisms.tex, source lines 14492--14492. Navigation pointer used in the counted span; the span contains the pointer only and the cited result is not reproduced inside it. Registered as a non-load-bearing navigation stub in the established form for Stacks material. No bibliography citation is involved.\par\medskip
\ifdefined\hypertarget\hypertarget{sourcepointer-4}{}\fi
\noindent\textit{Source reference.} Original section 0EKF: sources/more-morphisms.tex, source lines 14203--14203. Navigation pointer used in the counted span; the span contains the pointer only and the cited result is not reproduced inside it. Registered as a non-load-bearing navigation stub in the established form for Stacks material. No bibliography citation is involved.\par\medskip
\ifdefined\hypertarget\hypertarget{sourcepointer-5}{}\fi
\noindent\textit{Source reference.} Original lemma 0B5T: sources/coherent.tex, source lines 4329--4329. Navigation pointer used in the counted span; the span contains the pointer only and the cited result is not reproduced inside it. Registered as a non-load-bearing navigation stub in the established form for Stacks material. No bibliography citation is involved.\par\medskip
\ifdefined\hypertarget\hypertarget{sourcepointer-6}{}\fi
\noindent\textit{Source reference.} Original lemma 0887: sources/coherent.tex, source lines 6631--6631. Navigation pointer used in the counted span; the span contains the pointer only and the cited result is not reproduced inside it. Registered as a non-load-bearing navigation stub in the established form for Stacks material. No bibliography citation is involved.\par\medskip
\ifdefined\hypertarget\hypertarget{sourcepointer-7}{}\fi
\noindent\textit{Source reference.} Original lemma 0EIY: sources/algebraization.tex, source lines 5622--5622. Navigation pointer used in the counted span; the span contains the pointer only and the cited result is not reproduced inside it. Registered as a non-load-bearing navigation stub in the established form for Stacks material. No bibliography citation is involved.\par\medskip
\ifdefined\hypertarget\hypertarget{sourcepointer-8}{}\fi
\noindent\textit{Source reference.} Original lemma 0EJI: sources/algebraization.tex, source lines 7129--7129. Navigation pointer used in the counted span; the span contains the pointer only and the cited result is not reproduced inside it. Registered as a non-load-bearing navigation stub in the established form for Stacks material. No bibliography citation is involved.\par\medskip
\ifdefined\hypertarget\hypertarget{sourcepointer-9}{}\fi
\noindent\textit{Source reference.} Original equation 0EKH: sources/more-morphisms.tex, source lines 14457--14457. Navigation pointer used in the counted span; the span contains the pointer only and the cited result is not reproduced inside it. Registered as a non-load-bearing navigation stub in the established form for Stacks material. No bibliography citation is involved.\par\medskip
\ifdefined\hypertarget\hypertarget{sourcepointer-10}{}\fi
\noindent\textit{Source reference.} Original lemma 0EHM: sources/properties.tex, source lines 2413--2413. Navigation pointer used in the counted span; the span contains the pointer only and the cited result is not reproduced inside it. Registered as a non-load-bearing navigation stub in the established form for Stacks material. No bibliography citation is involved.\par\medskip
\ifdefined\hypertarget\hypertarget{sourcepointer-11}{}\fi
\noindent\textit{Source reference.} Original lemma 0EL4: sources/algebraization.tex, source lines 9221--9221. Navigation pointer used in the counted span; the span contains the pointer only and the cited result is not reproduced inside it. Registered as a non-load-bearing navigation stub in the established form for Stacks material. No bibliography citation is involved.\par\medskip
\ifdefined\hypertarget\hypertarget{sourcepointer-12}{}\fi
\noindent\textit{Source reference.} Original proposition 0BXF: sources/coherent.tex, source lines 4226--4226. Navigation pointer used in the counted span; the span contains the pointer only and the cited result is not reproduced inside it. Registered as a non-load-bearing navigation stub in the established form for Stacks material. No bibliography citation is involved.\par\medskip
\ifdefined\hypertarget\hypertarget{sourcepointer-13}{}\fi
\noindent\textit{Source reference.} Original lemma 0EIQ: sources/algebraization.tex, source lines 4721--4721. Navigation pointer used in the counted span; the span contains the pointer only and the cited result is not reproduced inside it. Registered as a non-load-bearing navigation stub in the established form for Stacks material. No bibliography citation is involved.\par\medskip
\noindent\textit{Reading revision r002.} Original author context is supplied below and editorial reference bindings are qualified by the chapter. Mathematical statement and proof are unchanged. Numbering is editorial; cited inputs remain in the complete fixed source.\par\medskip
\subsection*{Original author context (uncounted)}
\noindent
Let $k$ be a field. Let $P$ be a proper scheme over $k$.
Let $\mathcal{L}$ be an ample invertible $\mathcal{O}_P$-module.
Let $s \in \Gamma(P, \mathcal{L})$ be a section and let
$Q = Z(s)$ be the zero scheme, see
Divisors, Definition \ref{divisors-definition-zero-scheme-s}.
Let $\mathcal{I} \subset \mathcal{O}_P$ be the ideal sheaf of $Q$.
We will use $\textit{Coh}(P, \mathcal{I})$ to denote the category
of coherent formal modules introduced in
Cohomology of Schemes, Section \ref{coherent-section-coherent-formal}.

\subsection*{Original author statement and proof}
\begin{proposition}
\label{proposition-lefschetz-existence}
In the situation above let $(\mathcal{F}_n)$ be an object of
$\textit{Coh}(P, \mathcal{I})$. Assume for all $q \in Q$ and for
all primes $\mathfrak p \in \mathcal{O}_{P, q}^\wedge$,
$\mathfrak p \not \in V(\mathcal{I}_q^\wedge)$ we have
$$
\text{depth}((\mathcal{F}_q^\wedge)_\mathfrak p) +
\dim(\mathcal{O}_{P, q}^\wedge/\mathfrak p) +
\dim(\overline{\{q\}}) > 2
$$
Then $(\mathcal{F}_n)$ is the completion of a coherent
$\mathcal{O}_P$-module.
\end{proposition}

\begin{proof}
By Cohomology of Schemes, Lemma \ref{coherent-lemma-existence-easy}
to prove the lemma, we may replace $(\mathcal{F}_n)$ by an object
differing from it by $\mathcal{I}$-torsion (see below for more precision).
Let $T' = \{q \in Q \mid \dim(\overline{\{q\}}) = 0\}$
and $T = \{q \in Q \mid \dim(\overline{\{q\}}) \leq 1\}$.
The assumption in the proposition is exactly that
$Q \subset P$, $(\mathcal{F}_n)$, and $T' \subset T \subset Q$
satisfy the conditions of
Lemma \ref{lemma-improvement-formal-coherent-module-better}
with $d = 1$; besides trivial manipulations of inequalities, use that
$V(\mathfrak p) \cap V(\mathcal{I}^\wedge_y) = \{\mathfrak m^\wedge_y\}
\Leftrightarrow \dim(\mathcal{O}_{P, q}^\wedge/\mathfrak p) = 1$
as $\mathcal{I}_y^\wedge$ is generated by $1$ element.
Combining these two remarks, we may replace $(\mathcal{F}_n)$ by the
object $(\mathcal{H}_n)$ of $\textit{Coh}(P, \mathcal{I})$ found in
Lemma \ref{lemma-improvement-formal-coherent-module-better}.
Thus we may and do assume $(\mathcal{F}_n)$ is pro-isomorphic to
an inverse system $(\mathcal{F}_n'')$ of coherent $\mathcal{O}_P$-modules
such that $\text{depth}(\mathcal{F}''_{n, q}) + \dim(\overline{\{q\}}) \geq 2$
for all $q \in Q$.

\medskip\noindent
We will use More on Morphisms, Lemma \ref{more-morphisms-lemma-apply-proj-spec}
and we will use the notation used and results found
More on Morphisms, Section \ref{more-morphisms-section-proj-spec}
without further mention; this proof will not make sense without
at least understanding the statement of the lemma.
Observe that in our case
$A = \bigoplus_{m \geq 0} \Gamma(P, \mathcal{L}^{\otimes m})$
is a finite type $k$-algebra
all of whose graded parts are finite dimensional $k$-vector spaces, see
Cohomology of Schemes, Lemma \ref{coherent-lemma-coherent-proper-ample}.

\medskip\noindent
By Cohomology of Schemes, Lemma \ref{coherent-lemma-inverse-systems-pullback}
the pull back by $\pi|_U : U \to P$ is an object
$(\pi|_U^*\mathcal{F}_n)$ of $\textit{Coh}(U, f\mathcal{O}_U)$
which is pro-isomorphic to the inverse system
$(\pi|_U^*\mathcal{F}_n'')$ of coherent $\mathcal{O}_U$-modules.
We claim
$$
\text{depth}(\pi|_U^*\mathcal{F}''_{n, y}) + \delta_Z^Y(y) \geq 3
$$
for all $y \in U \cap Y$. Since all the points of $Z$ are closed, we
see that $\delta_Z^Y(y) \geq \dim(\overline{\{y\}})$ for all
$y \in U \cap Y$, see Lemma \ref{lemma-discussion}.
Let $q \in Q$ be the image of $y$. Since the morphism $\pi : U \to P$ is
smooth of relative dimension $1$ we see that either $y$ is a closed point
of a fibre of $\pi$ or a generic point.
Thus we see that
$$
\text{depth}(\pi^*\mathcal{F}''_{n, y}) + \delta_Z^Y(y)
\geq
\text{depth}(\pi^*\mathcal{F}''_{n, y}) + \dim(\overline{\{y\}}) =
\text{depth}(\mathcal{F}''_{n, q}) + \dim(\overline{\{q\}}) + 1
$$
because either the depth goes up by $1$ or the dimension.
This proves the claim.

\medskip\noindent
By Lemma \ref{lemma-cd-1-canonical} we conclude that
$(\pi|_U^*\mathcal{F}_n)$ canonically extends to $X$.
Observe that
$$
M_n = \Gamma(U, \pi|_U^*\mathcal{F}_n) =
\bigoplus\nolimits_{m \in \mathbf{Z}}
\Gamma(P, \mathcal{F}_n \otimes_{\mathcal{O}_P} \mathcal{L}^{\otimes m})
$$
is canonically a graded $A$-module, see
More on Morphisms, Equation (\ref{more-morphisms-equation-cohomology-torsor}).
By Properties, Lemma \ref{properties-lemma-quasi-coherent-quasi-affine}
we have $\pi|_U^*\mathcal{F}_n = \widetilde{M_n}|_U$.
Thus we may apply Lemma \ref{lemma-Gm-equivariant-extend-canonically}
to find a finite graded $A$-module $N$ such that
$(M_n)$ and $(N/I^nN)$ are pro-isomorphic in the category
of graded $A$-modules modulo $A_+$-torsion modules.
Let $\mathcal{F}$ be the coherent $\mathcal{O}_P$-module
associated to $N$, see
Cohomology of Schemes, Proposition
\ref{coherent-proposition-coherent-modules-on-proj-general}.
The same proposition tells us that $(\mathcal{F}/\mathcal{I}^n\mathcal{F})$
is pro-isomorphic to $(\mathcal{F}_n)$.
Since both are objects of $\textit{Coh}(P, \mathcal{I})$
we win by Lemma \ref{lemma-recognize-formal-coherent-modules}.
\end{proof}

\end{document}
